% Transcription — fonds Alexandre Grothendieck, shelfmark 67, batch 2 (pages 21-40).
% Established from the facsimile archives/batches/67/batch-02.pdf (a mirror of
% https://grothendieck.umontpellier.fr/67.pdf), per the skill
% transcribe-grothendieck. The batch file carries no cover sheet: PDF page N
% is page 20+N of the folder, confirmed by the Montpellier footer printed
% under each leaf. \page{} numbers that position, as folders 1 and 67/4 do:
% the archivists' pencil numbers bottom-left read 21 on page 21, then run
% three ahead (page 22 is pencilled 25, page 40 is pencilled 43).
%
% Pass: Opus 5.5 (claude-opus-5-5), 2026-09-22 — first pass, unchecked against the pages by a human.
%
% What the batch is.
%   21     A second, fainter copy (carbon) of the last typed leaf of the
%          letter to Lipman Bers (batch 1, page 20), paginated « 5 » at the
%          top, with the same two handwritten insertions. Transcribed, since
%          it carries the mathematics; the letter itself is not repeated.
%   22     Blank. Absent.
%   23     Cover in his hand, top right, on the back of a computer listing:
%          « Métriques canoniques sur les surfaces conformes ». No \page;
%          its words open the second section.
%   24     A small ruled leaf: a list a), b), c), d), d') of equivalent data
%          on a conformal sphere (spherical metric, great circle, compact
%          torus of Aut^+ and one of its orbits, pair of points).
%   25     Sketch sheet filmed upside down: u_+, u_-, an exact sequence
%          0 -> Π -> V -> E -> 0 and H^1(γ,E) ≃ H^2(γ,Π) ≃ Π^γ/(1+σ)Π,
%          z |-> exp 2iπz.
%   26-40  One run in ink on tractor-feed paper, rectos paginated 1) to 15)
%          by him (page 26 = his 1, …, page 40 = his 15; the numbers are
%          noted once per page). Hyperbolic surfaces (Th: five equivalent
%          conditions, Déf, Cor. 1 and 2); the non-hyperbolic list; then
%          « Question des métriques riemanniennes canoniques sur une surface
%          conforme »: (1) hyperbolic case; 2) elliptic curves, 2') C* and
%          the real projective plane minus a point, 2'') E_C; 3) the real
%          projective plane, 3') P^1_C where no canonical metric exists
%          without extra structure H ⊂ Aut_conf(X) compact; the analysis of
%          H through H^+, the maximal torus T, the extension N of Z/2Z by T,
%          H^2(γ,T), inversions against antipodisms (pp. 31-35); « Prop. »
%          en résumé (p. 36); « Récapitulation » (p. 37) with R(X), Q_X;
%          curvature of Q_X (p. 38); « Cas des surfaces conformes à bord »
%          by the double X^or with its anti-involution (pp. 39-40). The run
%          stops mid-sentence at the foot of page 40 (« la deuxième »); its
%          continuation, if any, is in batch 3.
%
% Notation fixed for the batch. X^{\mathrm{or}} for the superscript he
% writes after X for the orientation double (read « or », the r looking like
% a 2 on page 39): the holomorphic surface without boundary carrying the
% anti-involution σ, with X = X^or/σ. X^σ, X^H, X^T are fixed-point sets.
% \mathbb{U} for his double-struck U (complex numbers of modulus 1); G =
% Aut_conf(X), G^+ = Aut^+(X), H^+ = H ∩ G^+; N^- the orientation-reversing
% elements of N; γ = Z/2Z acting; Herm^{++}(V) the positive definite
% hermitian forms; R(X) the normalised spherical metrics; Q_X, Q_{X,H} the
% canonical metric. E^1_C is his affine line, G_{m,C} the multiplicative
% group, E_C an elliptic curve. His « ssi », « i.e. », « cte », « OPS »
% (on peut supposer), « tt », « comp. », « corr. » kept as he writes them.
%
% Legibility. A fast hand throughout; the formulas carry, the connecting
% prose often only by the argument. The oblique margins of pages 36 and 38
% were turned upright before reading. The margin of page 38 is the least
% certain passage of the batch. No date appears on these twenty leaves.
%
% The dating is the inventory's own for the folder, copied verbatim. Its
% group, [66-89] « Autour de l'enseignement - Thèmes liés (cours et
% directions de recherches) », is dated 1964-1984.
\documentclass[11pt,a4paper]{article}
% Le préambule commun à toutes les transcriptions du fonds Grothendieck.
%
% Il tient en une idée : **l'incertitude se compose, elle ne se résout pas.**
% Une transcription qui lisse un mot douteux a détruit la seule chose pour
% laquelle on la faisait — un lecteur doit pouvoir savoir, à chaque endroit,
% ce qui était sur le feuillet et ce qui a été ajouté. Les six macros qui
% suivent sont donc l'appareil critique, et non de la décoration.
%
% Les mêmes commandes sont reconnues par `scripts/render.mjs`, qui produit la
% vue de lecture du site. Une macro ajoutée ici doit l'être aussi là-bas,
% faute de quoi le rendu échouera bruyamment — ce qui est voulu : un rendu
% silencieusement faux est pire qu'un rendu absent.

\ProvidesPackage{grothendieck}[2026/08/08 Transcriptions du fonds Grothendieck]

\usepackage{amsmath,amssymb,amsthm}
\usepackage{mathrsfs}
\usepackage{stmaryrd}
% Les diagrammes commutatifs sont partout dans ce fonds : c'est la langue
% même de la géométrie algébrique de Grothendieck, et une transcription qui
% les rendrait en prose ne transcrirait rien.
\usepackage{tikz-cd}
% Français par défaut — c'est la langue des feuillets — mais l'anglais est
% chargé aussi : la traduction et le résumé sont des documents anglais, et la
% typographie française (espaces avant les deux-points) y serait une faute.
% Ces documents basculent par \selectlanguage{english} en tête de corps.
\usepackage[english,french]{babel}
\usepackage{fontspec}
\usepackage{geometry}
\geometry{a4paper,margin=25mm,marginparwidth=28mm}
\usepackage{marginnote}
\usepackage{xcolor}
% ulem, not soul. `soul`'s \st and \ul are built on a character-by-character
% scan that XeLaTeX's font machinery breaks: the document fails with an
% undefined control sequence rather than degrading. ulem does the same two jobs
% and is Unicode-engine safe. `normalem` stops it hijacking \emph, which the
% transcriptions use for Grothendieck's own emphasis.
\usepackage[normalem]{ulem}
\usepackage{hyperref}
\hypersetup{colorlinks=true,linkcolor=black,urlcolor=black,pdfborder={0 0 0}}

% --- Métadonnées du lot -----------------------------------------------------
% Reprises telles quelles par le script de rendu, qui les affiche en tête de
% la vue de lecture. Elles fixent ce que le fichier prétend couvrir : sans
% elles, une transcription flotte sans point d'attache dans un fonds de seize
% mille feuillets.

\newcommand{\folder}[1]{\gdef\grothFolder{#1}}
\newcommand{\batch}[1]{\gdef\grothBatch{#1}}
\newcommand{\leaves}[2]{\gdef\grothFirst{#1}\gdef\grothLast{#2}}
\newcommand{\foldertitle}[1]{\gdef\grothFoldertitle{#1}}
\gdef\grothFolder{?}\gdef\grothBatch{?}\gdef\grothFirst{?}\gdef\grothLast{?}\gdef\grothFoldertitle{}

% --- Le résumé --------------------------------------------------------------
%
% Il ouvre la lecture modernisée, sous le titre « Résumé », et il est composé
% en petit. Ce n'est pas un ornement : c'est le seul passage du document qui
% ne soit pas des mathématiques, et un lecteur doit pouvoir le reconnaître
% avant de l'avoir lu — puis le sauter s'il connaît déjà le sujet, ou s'y
% arrêter s'il ne le connaît pas. Le filet qui le suit marque la reprise du
% corps.

\newenvironment{resume}
  {\small\setlength{\parskip}{0.45em}}
  {\par\medskip\hrule\medskip}

% --- Mots-clés ---------------------------------------------------------------
%
% En anglais, en clôture du résumé de la lecture modernisée : le vocabulaire
% moderne sous lequel chercher ce que les feuillets construisent. Le manifeste
% du site (`npm run manifest`) les extrait pour étiqueter la cote entière —
% cette ligne est la source unique des tags, il n'y a pas de fichier à côté.
\newcommand{\keywords}[1]{\par\smallskip\noindent
  {\footnotesize\textsc{Keywords} --- \textit{#1}}\par}

% --- L'appareil critique ----------------------------------------------------

% Le numéro de feuillet, en marge. C'est l'ancre qui permet de régler un
% désaccord sur une formule en regardant *un* feuillet plutôt qu'une chemise
% de six cents. Le site s'en sert aussi pour tourner les pages du fac-similé
% quand on fait défiler la transcription.
\newcommand{\leaf}[1]{%
  \marginnote{\footnotesize\color{gray!70}\textsf{#1}}%
  \label{leaf:#1}\ignorespaces}

% Illisible : on ne devine pas. Un mot inventé qui se lit comme les autres est
% le pire résultat possible.
\newcommand{\ill}{\textcolor{red!60!black}{[\,\ldots\,]}}

% Lecture douteuse : le mot est donné, et signalé comme incertain.
\newcommand{\uncertain}[1]{\uline{#1}}

% Ajout de l'éditeur — une lettre manquante, un mot sous-entendu.
\newcommand{\add}[1]{\textcolor{blue!50!black}{[#1]}}

% Biffé par Grothendieck lui-même. Ce qu'il a rayé fait partie du document :
% une rature dit souvent où la pensée a changé de direction.
\newcommand{\struck}[1]{\sout{#1}}

% Note du transcripteur, sur l'état matériel du feuillet.
\newcommand{\note}[1]{\par\smallskip{\small\color{gray!60!black}\textit{[#1]}}\par\smallskip}

% Note marginale de Grothendieck, distincte du corps du texte.
\newcommand{\marginal}[1]{\par\smallskip\noindent\hspace{1em}%
  {\small\color{gray!60!black}#1}\par\smallskip}

% --- Titre ------------------------------------------------------------------

\AtBeginDocument{%
  \begin{center}
    {\large\bfseries Fonds Alexandre Grothendieck --- cote n\textsuperscript{o}~\grothFolder}\\[2pt]
    {\itshape\grothFoldertitle}\\[4pt]
    {\small feuillets \grothFirst--\grothLast\ (lot \grothBatch)}\\[2pt]
    {\footnotesize\color{gray!60!black}Transcription établie d'après le fac-similé numérisé
     par l'Université de Montpellier.\\
     Les lectures incertaines sont soulignées, les illisibles notées [\,\ldots\,],
     les ajouts entre crochets.}
  \end{center}
  \vspace{1em}}

\endinput


\folder{67}
\batch{2}
\pages{21}{40}
\dating{[à partir de 1977]-1984}
\watermark{Édition de démonstration}
\foldertitle{Chirurgie des surfaces conformes : notes manuscrites (s.d., 1983-1984), lettres (1984).}

\begin{document}

\section{Lettre à Lipman Bers (second exemplaire du dernier feuillet)}

\page{21}
\note{second exemplaire, plus pâle, du dernier feuillet dactylographié de la
lettre à Lipman Bers (lot 1, page 20), paginé « 5 » en haut ; le texte est
le même, avec les deux mêmes ajouts à la main}

hold on more specific properties. As an example, starting with a compact
conformal surface with boundary $X$ (a pant, say), of hyperbolic type, and
removing an (open) “collar” around the boundary, we get another surface with
boundary $X'$ — what about the relation \struck{of} between the two
corresponding metrics ? Assuming the model for pants above is correct, it
would be nice to have an explicit expression of the metric of a pant in terms
of the parameters $\lambda_{i}$.

With my thanks for your attention, and for whatever comment you will care to
make, very sincerely ypur's
\note{« compact » et « between » sont ajoutés à la main au-dessus de la ligne}

\section{Métriques canoniques sur les surfaces conformes}
\note{titre de sa main, en haut à droite du feuillet 23, par ailleurs vierge
(le dos d'un listing d'ordinateur), qui ne reçoit pas de numéro de page}

\page{24}
\note{petit feuillet réglé ; une liste de données équivalentes sur une sphère
conforme $X$, $G = \mathrm{Aut}^{+}(X)$}

a) deux \uncertain{réalifications} de types \ill{}, compatibles, i.e.\ une
paire compatible d'antipodisme et \uncertain{inversion}

b) Métriques sphériques $+$ grand cercle

c) cercle $+$ structure torique compatible \uncertain{dessus}, i.e.\
antiinvolution \struck{d\ill{}} \add{\uncertain{avec} \uncertain{pts}
\uncertain{fixes}}, et paire de points conjugués par celle-ci

d) Tore compact $T \subset G = \mathrm{Aut}^{+}(X)$, $+$ orbite de
$X \smallsetminus X^{T}$ sous $T$ (i.e.\ orbites non réduites à un point)

d') paire de pts $\{z, z'\} \subset X$, $+$ orbite de
$X \smallsetminus \{z, z'\}$ du tore compact unique $T \subset G$ qui fixe
$z, z'$.

\page{25}
\note{feuille de brouillon : formules isolées, sans texte, relevées dans
l'ordre de la page}

\[
  u_{+} \qquad \tfrac{1}{2}(u_{+} + u_{-})
\]
\[
  \underbrace{\left(x + \tfrac{1}{2} y\right)}_{m} u_{+} +
  \underbrace{\tfrac{1}{2} y}_{n} u_{-}
\]
\[
  0 \to \Pi \to V \to E \to 0
\]
\[
  H^{1}(\gamma, E) \xrightarrow{\ \sim\ } H^{2}(\gamma, \Pi) \simeq
  \Pi^{\gamma} / (1 + \sigma)(\Pi)
\]
\[
  z \longmapsto \exp 2i\pi z
\]

\page{26}
\note{pagination de l'auteur : 1, cerclé}

$X$ surface \struck{conforme} \add{hol 0-connexe} (pas néc.\
cpte)
\note{dans ce lot, ce qui est marqué comme ajout est, sauf mention contraire,
de sa main, porté en interligne au-dessus du mot biffé ou de la ligne}
\marginal{en haut à droite, séparé par un trait oblique : \emph{NB} Ici
“surface” sous-entend “séparée”}

$\widetilde{X}$ rev.\ universel

\emph{Th} Conditions équivalentes

(i) $X$ n'est pas isom.\ à une des trois surfaces
$\mathbb{P}^{1}_{\mathbb{C}}$,
$\mathbb{E}^{1}_{\mathbb{C}} = \mathbb{P}^{1}_{\mathbb{C}} \smallsetminus
\{\infty\}$, $\mathbb{G}_{m\,\mathbb{C}} = \mathbb{P}^{1}_{\mathbb{C}}
\smallsetminus \{\infty, 0\}$, ni à une courbe elliptique complexe

(ii) $\widetilde{X}$ est isomorphe à $H$ (demi-plan de Poincaré)

(iii) Toute application holomorphe $\mathbb{C} \to X$ est constante

(iv) [Si $X$ est un ouvert d'une droite projective complexe (resp.\ d'une
droite affine complexe)]
$\operatorname{card}(X' \smallsetminus X) \geq 3$ (resp.\
$\operatorname{card}(X' \smallsetminus X) \geq 2$).
\note{dans les deux cardinaux, un premier $X$ est biffé et récrit $X'$ ; $X'$
désigne la droite ambiante}

\struck{(iv') [Si $X$ est un ouvert de la surface]}

(v) [Si $X = \widehat{X} \smallsetminus S$, $\widehat{X}$ compacte
\add{de genre $g$}, $S$ finie \add{de card $\nu$}]
$(2g + \nu \geq 3$ i.e.\ $(g, \nu) \notin \{(0,0), (0,1), (0,2), (1,0)\}$
i.e.\ $\chi(X) < 0)$
\note{sous la ligne, biffé : « holomorphe \ill{} conforme » ; « genre $g$ »
et « card $\nu$ » sont en interligne ; le $2$ de $2g$ est écrit sur un $3$}

\emph{Déf} On dit alors que $X$ est hyperbolique.

\emph{Corollaire 1} Un ouvert 0-connexe d'une surface hol hyperbolique est
hyperbolique.

\emph{Corollaire 2} Soit $X$ surface hol.\ \emph{à bord}, \struck{\ill{}}
connexe, avec $\partial X \neq \emptyset$. Alors $\operatorname{Int} X$ est
hyperbolique

\page{27}
\note{pagination de l'auteur : 2, cerclé. La phrase qui introduit la liste
n'est pas sur la page précédente}

les surfaces suivantes :

a) le plan projectif réel ($\simeq$ Sphère conforme/antipodie) et le plan
projectif réel privé d'un point (dont les rev.\ universels sont resp.\
$\simeq \mathbb{P}^{1}_{\mathbb{C}}$ et $\simeq \mathbb{G}_{m\,\mathbb{C}}$).

b) Les surfaces conformes associées aux courbes \struck{algébriques}
\add{lisses géom.\ connexes} réelles \add{de genre 1} $E$ telles que
$E(\mathbb{R}) = \emptyset$, on prend \struck{$E(\mathbb{C})$}
\[
  X = E(\mathbb{C}) / \text{conj.\ complexe} .
\]
On trouve que les $E$ en question sont, à isom.\ près, les $E_{a}$,
$a \in \mathbb{R}^{*+}$\,…

\note{un trait horizontal barre la page}

\emph{Question des \add{les} métriques riemanniennes canoniques sur une
surface conforme} (sans bord pour commencer)

(1) Cas hyperbolique. Il y a une métrique canonique, à courbure constante de
valeur $-1$, déduite par descente du revêtement universel $\widetilde{X}$.
\note{sous cette ligne, soulignés : « Métriques hyperboliques », comme un
titre ou un renvoi}

\page{28}
\note{pagination de l'auteur : 3}

déterminée par la condition que l'aire totale de $X$ soit $= 1$.
\note{le début du cas (2) manque : la page précédente passe du cas (1) à
celle-ci sans transition}

On trouve une \emph{métrique plate} caractérisée par la condition d'être
invariante par translations, et d'aire totale égale à $1$ — donc stable par
tt autom.\ conformes (pas seulement holomorphes). Il s'ensuit que dans le cas
non orienté, la métrique canonique de $X^{\mathrm{or}}$ (qui est une courbe
elliptique $/\mathbb{C}$) se descend à \uncertain{$X$}.
\note{le dernier signe porte le même exposant que $X^{\mathrm{or}}$ ; le sens
demande $X$}

2') Cas $X \simeq \mathbb{C}^{*}$, ou \add{$X \simeq$} plan projectif réel
privé d'un point [i.e.\ $X$ non orientable et
$X^{\mathrm{or}} \simeq \mathbb{C}^{*}$]. Le cas $X \simeq \mathbb{C}^{*}$,
comme dans 2, \struck{plus intrinsèquement $X \simeq$} par descente à partir
des rev.\ universels on trouve une métrique canonique mod.\ multiplication
par un scalaire $> 0$, \emph{métrique plate} également. \struck{\ill{}} les
autom.\ de $\mathbb{C}^{*}$ \struck{\ill{}} \add{\uncertain{hol.}} sont de la
forme $z \to az$, soit $z \to a/z$, qui conservent la métrique [car sur
$\widetilde{X} \simeq \mathbb{C}$ ils correspondent à des $z \to z + \alpha$,
$z \mapsto -z + \alpha$], donc \uncertain{en divisant} la métrique sur
$\mathbb{C}^{*}$, \struck{de façon}

\page{29}
\note{pagination de l'auteur : 4}

\ill{} en la normalisant \add{p.\ ex.} de telle façon que la longueur du
cercle $|z| = 1$ soit $1$, on la définit sur toutes les
$X \simeq \mathbb{C}^{*}$ du même fait. D'ailleurs l'antiautomorphisme
$z \mapsto \bar{z}$ \struck{\ill{}} conserve la \uncertain{mesure}, donc si
$X \simeq \mathbb{C}^{*}$, tt autom.\ conforme respecte la \uncertain{mesure}
canonique. Il en résulte que dans le cas non orienté
$X \simeq$ plan projectif réel moins 1 pt, on peut descendre à partir de
$X^{\mathrm{or}}$ et trouver encore une métrique canonique plate, stable par
tous les automorphismes conformes.
\marginal{en marge droite : C'est la métrique déduite de \ill{}
\uncertain{hermitienne} canonique $z\bar{z}$ sur $\mathbb{C}$, via
$z \mapsto \exp 2i\pi z$, de $\mathbb{C}$ sur $\mathbb{C}^{*}$.}

2'') Cas $X \simeq \mathbb{E}_{\mathbb{C}}$. Les métriques riemanniennes sur
une courbe \struck{\ill{}} \uncertain{elliptique} complexe, correspondent aux
structures hermitiennes \add{déf.\ pos.} sur $L$, (celles-ci sont bien
définies à homothétie près) et sont invariantes par \uncertain{translation}
\struck{\ill{}} \add{réelle} \uncertain{de} $L$, mais seulement par le
sous-groupe engendré par les translations, et les \struck{autres}
homothéties par des $u \in \mathbb{U} = \{z \in \mathbb{C} \mid |z| = 1\}$ ;
elles sont quasi-invariantes par autom.\ quelconques (lesquels sont
\emph{affines}). Donc si $X \simeq \mathbb{E}_{\mathbb{C}}$, on peut définir
une \struck{\ill{}} métrique

\page{30}
\note{pagination de l'auteur : 5}

canonique stable par translations etc., mais une donnée de telles mod
homothétie par $c \in \mathbb{R}^{+*}$ — ces métriques sont caractérisées
par la condition d'être invariantes par translations [\emph{NB} $X$ a
canoniquement une structure de droite affine sur $\mathbb{C}$]. Ces
métriques \struck{\ill{}} ne dépendent que de la structure conforme
sous-jacente.

3) \emph{Cas \struck{\ill{}} $X \simeq$ \struck{sphère} plan projectif réel}.
Dans ce \struck{premier} cas, $X^{\mathrm{or}}$ est une droite projective
complexe munie d'un antiautomorphisme involutif sans pt fixe
(antipodisme), ou ce qui revient au même, d'une métrique “euclidienne
sphérique” compatible avec sa structure conforme [celle-ci, \struck{\ill{}}
\add{\uncertain{sera} \uncertain{normalisée} par} la condition que l'aire de
la surface soit $= 1$]. Étant invariante par l'antipodisme, cette métrique se
descend à $X$.

3') Cas $X \simeq \mathbb{P}^{1}_{\mathbb{C}}$. Il n'y a pas de métriques
riemanniennes sur $X$, invariantes [ou quasi-invariantes] par tt
automorphisme — donc \uncertain{sans} structure \uncertain{supplémentaire}
sur $X$

\page{31}
\note{pagination de l'auteur : 6}

il n'est pas possible d'y introduire une métrique riemannienne canonique,
même modulo cte. Considérons cependant la donnée d'une structure
supplémentaire du type
\[
  H \subset \mathrm{Aut}_{\mathrm{conf}}(X) = G
\]
où $H$ est un sous-groupe \emph{compact} de
$\mathrm{Aut}_{\mathrm{conf}}(X)$ (condition \emph{nécessaire} pour que $H$,
supposé fermé, invarie une métrique riemannienne sphérique comp.\ avec la
str.\ conforme). Soit $H^{+} = H \cap \mathrm{Aut}^{+}(X)$ (sous-groupe des
autom.\ qui conservent l'orientation). OPS $X = \mathbb{P}(V)$, $V$ vectoriel
complexe de rg 2, \struck{\ill{}} soit $\widetilde{H}$ l'image inverse de
$H^{+}$ dans $\mathrm{SL}(V)$, c'est donc un sous-groupe cpt, qui invarie
nécessairement une forme hermitienne déf.\ pos.\ sur $V$, d'où la métrique
sphérique couv.\ sur $X$. (normalisée par la condition
$\int_{X} dx = 1$) Dans le cas où \uncertain{la} \struck{métrique}
\struck{autre métrique} hermitienne invariante est unique mod cte, on trouve
donc une métrique \add{sphérique normalisée} canonique sur $X$, définie par
la \add{seule} donnée de $H^{+}$. Cette métrique sera $=$ invariante par
$\mathrm{Norm}_{G}(H^{+})$, et a fortiori par $H$, et est aussi
caractérisée par cette condition.
\note{« (normalisée par la condition $\int_{X} dx = 1$) » est en interligne,
sous une accolade ; dans $\mathrm{Norm}_{G}(H^{+})$, un mot biffé suit
$H^{+}$}

\page{32}
\note{pagination de l'auteur : 7}

Nous allons maintenant examiner de plus près le cas où \struck{\ill{}}
$\mathrm{Herm}^{++}(V)^{H^{+}}$ n'est pas réduit à une demi-droite. On
vérifie aisément que cela revient à dire que $H^{+}$ \struck{\ill{}} fixe
deux pts distincts $z, z'$ de $X$, donc est contenu dans le sous-groupe de
$G^{+}$ \struck{défini} formé des élts qui fixent $z, z'$, lequel est isom.\
à $\mathbb{C}^{*}$, \struck{\ill{}} \add{\uncertain{dont} \uncertain{l'unique}}
sous-groupe cpt maximal \add{$T$ isom.\ à $\mathbb{U}$}
[alors $\{z, z'\}$ dans $T$ est uniquement déterminé]
\struck{Dans ce cas. Les métriques invariantes par}
\marginal{en marge droite, le long d'un trait vertical : à vérifier}

Si $H^{+} \neq \{1\}$, les métriques invariantes par $H^{+}$ sont aussi
celles invariantes par $T \supset H^{+}$, et elles sont en corr.\ 1-1 avec
les \struck{trajectoires} \add{orbites} de $T$ dans
$X^{*} = X \smallsetminus \{z, z'\}$, [\emph{NB} $T$ opère librement sur
$X^{*}$].

La relation se détermine comme suit : \uncertain{à} une métrique invariante
par $T$, l'\struck{équateur} \uncertain{grand cercle} \uncertain{équidistant}
des pts $z z'$, lequel est une \struck{trajectoire} \add{orbite de $X^{*}$}
sous $T$. En particulier, les orbites sous $H^{+}$ (\uncertain{contenues}
dans une unique orbite de $T$) définissent une telle métrique.

Si $H = H^{+}$, il n'y a \struck{\ill{}} rien à ajouter. Dans le

\page{33}
\note{pagination de l'auteur : 8}

cas contraire, distinguons deux cas.

a) $H^{+} = \{1\}$, donc $H = \{1, \sigma\}$, où $\sigma$ est une
anti-involution.

\quad 1) Si \struck{\ill{}} $X^{\sigma} = \emptyset$, $\sigma$ fixe une
métrique sphérique normalisée unique, OK.

\quad 2) Si $X^{\sigma} \neq \emptyset$ i.e.\ $\sigma$ une inversion, alors
le choix d'une métrique sphérique comp.\ avec $\sigma$ revient à celui d'une
paire de pts de $X$ conjugués sous $\sigma$. En particulier, le choix d'un
point $x$ de $X$ tel que $\sigma x \neq x$ définit une telle métrique
sphérique

b) $H^{+} \neq \{1\}$. Considérons le \add{\uncertain{unique}} tore compact
maximal $T$ qui contient $H^{+}$, et $T.H = N$, ext.\ de
\struck{\ill{}} $\mathbb{Z}/2\mathbb{Z}$ par $T$. Les élts $\sigma$ de
$N^{-}$ fixent $z, z'$, ou les intervertissent, suivant que l'élt $-1$ de
$N/T$ opère trivialement sur $T$, ou par la symétrie $t \mapsto t^{-1}$.
\marginal{en marge droite, le long d'un trait vertical qui couvre ces lignes :
c'est l'inverse !}
Dans le 1er cas, l'ext.\ $N$ de $\overset{\gamma}{\mathbb{Z}/2\mathbb{Z}}$
par $T$ est triviale
\[
  (H^{2}(\gamma, T) \simeq T^{\gamma} / N_{\gamma}(T) = T / T^{2} \simeq 0)
\]
dans le deuxième cas on trouve
\[
  H^{2}(\gamma, T) \simeq T^{\gamma} / N_{\gamma}(T) \simeq
  \underbrace{{}_{2}T}_{\mathbb{Z}/2\mathbb{Z}} / 0 \simeq
  \mathbb{Z}/2\mathbb{Z} .
\]

1) $N$ ext.\ \emph{centrale} de \struck{$\gamma$}
\add{$\mathbb{Z}/2\mathbb{Z}$} par $T$, donc
$N \simeq T \times \underset{\{1, \sigma\}}{\gamma}$ (le direct)
engendré par $T$ et un $\sigma$ \struck{antiinvolution} qui
\struck{fixe $z, z'$} intervertit $z, z'$. \struck{On voit que}
\struck{\ill{}} celle-ci est une inversion, \struck{\uncertain{qui}}
\note{en bas à droite, un croquis : une sphère, $z$ au pôle nord, $z'$ au
pôle sud, quelques méridiens et un point marqué $x$}

\page{34}
\note{pagination de l'auteur : 9}

définit donc (via la paire $\{z, z'\}$) une métrique sphérique, ce qui sera
OK. Une fois \uncertain{établi} \add{explicité} que $\sigma$ n'est pas
unique, (il y a une ambiguïté d'ordre 2, on peut le remplacer par
$\sigma' = \varepsilon \circ \sigma$, où $\varepsilon$ est l'élément d'ordre
2 de $T$. On trouve en fait que $\{\sigma, \sigma'\}$ est formé d'une
inversion et d'un antipodisme.

\emph{Dém.} OPS $X =$ sphère de Riemann, $\{z, z'\} = \{\infty, 0\}$, donc
\[
  T = \{ z \mapsto az \mid a \in \mathbb{U} \text{ i.e. } a \in \mathbb{C},\
  |a| = 1 \text{ i.e. } a\bar{a} = 1 \}
\]
Les \struck{\ill{}} autom.\ conformes anti-hol.\ de $\mathbb{C}^{*}$ sont
\[
  z \longmapsto a\bar{z} \quad (a \in \mathbb{C}^{*}) \quad
  \text{ceux qui fixent } 0, \infty
\]
\[
  z \longmapsto b/\bar{z} \quad (b \in \mathbb{C}^{*}) \quad
  \text{— intervertissant } 0, \infty
\]
On voit tout de suite que ce sont les \emph{seconds} qui normalisent $T$.
\struck{Un tel} Si
\[
  \sigma(z) = b/\bar{z} , \quad \text{on a} \quad
  \sigma^{2}(z) = \tfrac{b}{\bar{b}}\, z ,
\]
donc $\sigma$ involutif ssi $b/\bar{b} = 1$ i.e.\ $b \in \mathbb{R}^{*}$.
Les pts fixes de $\sigma$ sont donnés par
\[
  z\bar{z} = b
\]
\struck{\ill{}} donc $\sigma$ est une inversion resp.\ un antipodisme,
suivant que \struck{$b \in \mathbb{R}$} $b > 0$ ou $b < 0$.

Si $\sigma(z) = b/\bar{z}$, on a $\sigma'(z) = -b/\bar{z}$, donc si
\struck{\ill{}} $\sigma$ est une inversion, $\sigma'$ est un antipodisme, et
vice versa. OPS les notations choisies de telle façon que $b > 0$, i.e.\
$\sigma$ inversion, donc $\sigma$ définit une métrique sphérique
\uncertain{dont} \uncertain{un} grand cercle \uncertain{est} défini par
$C = \{z \in \mathbb{C} \mid z\bar{z} = b\}$

\page{35}
\note{pagination de l'auteur : 10}

\struck{\ill{}} avec les pôles $0, \infty$. D'autre part, l'antipodisme
$\sigma' z = -b/\bar{z}$ \struck{transforme} \add{intervertit} $0$ et
$\infty$, et \struck{transforme $z \in C$} $C$ en lui-même, car
$z \in C$ i.e.\ $z\bar{z} = b$ implique que $\sigma' z \in C$ i.e.\
\[
  \underbrace{\underbrace{(\sigma' z)}_{(-b/\bar{z})}\,
  \underbrace{\overline{(\sigma' z)}}_{(-b/z)}}_{b^{2}/z\bar{z}} = b
\]
Ainsi, $\sigma'$ et $(\sigma, \{0, \infty\})$ définissent la même
structure riemannienne sphérique. On est content !

2) $N$ \struck{est} \emph{non centrale} de $\gamma$ par $T$. Cela signifie
que si $\sigma \in N \smallsetminus T$, on a
\[
  \sigma z = a\bar{z} , \quad (a \in \mathbb{C}^{*}) .
\]
\struck{Donc} On en conclut
\[
  \sigma^{2} z = a(\overline{a z}) = a\bar{a}\, z
\]
et la condition $\sigma^{2} \in T$ signifie $a\bar{a} = 1$ i.e.\
$a \in \mathbb{U}$, qui implique aussi \struck{\ill{}} $\sigma^{2} = 1$ i.e.\
l'extension $N$ est \uncertain{scindée}.
\note{on lit $\alpha^{2} = 1$ ; le calcul qui précède donne
$\sigma^{2} = 1$}

Le produit $u_{\theta}\sigma$ ($\theta \in \mathbb{U}$) est
$u_{\theta}\sigma(z) = (\theta a)\bar{z}$ de la même forme. Soit
$\alpha \in \mathbb{U}$ tel que $a = \alpha^{2} = \alpha/\bar{\alpha}$ donc
\[
  z \in \mathbb{C}^{\sigma} \iff z = \alpha \frac{\bar{z}}{\bar{\alpha}}
  \iff z/\alpha = \overline{(z/\alpha)} \quad \text{i.e.} \quad
  z \in \alpha \mathbb{R}
\]
\marginal{en marge droite, séparé par un trait : Ici $N$ est déterminé en
termes de $T$, i.e.\ en termes de $\{z, z'\}$, comme $T \cup N^{-}$, où
$N^{-}$ est formé de \emph{toutes} les anti-involutions de $X$ qui fixent
$\{z, z'\}$. Donc il n'est pas question, \ill{} \ill{}, d'extraire une
métrique riem.\ sphérique}

\page{36}
\note{pagination de l'auteur : 11}

En résumé et sauf erreur

\emph{Prop.} Soit $X$ une sphère conforme \struck{\uncertain{marquée}},
$H \subset \mathrm{Aut}_{\mathrm{conf}}(X) \overset{\mathrm{def}}{=} G$ un
sous-groupe \emph{compact} \struck{de}. Alors

a) Il existe une métrique riemannienne sphérique normalisée sur $X$ comp.\
avec sa structure conforme \uncertain{invariante} \uncertain{par} $H$

b) Celle-ci est unique ssi $X^{H} = \emptyset$ \struck{il n'existe pas de}

\struck{c)} c) Supposons \struck{\ill{}} $X^{H} \neq \emptyset$,
$H \neq \{1\}$. Alors on a \struck{\ill{}} \add{ou bien}
$\operatorname{card} X^{H} = 2$, ou bien $H = \{1, \sigma\}$ où $\sigma$ est
une inversion.

\quad 1) Si $\operatorname{card} X^{H} = 2$, \struck{\ill{}}
$H^{+} = H \cap G^{+}$, $G^{+} \ldotp X^{H} = X^{H^{+}}$, et $H^{+}$ est
contenu dans un unique tore compact $T$, et \struck{les métriques sphériques}
les métriques sphériques riemanniennes normalisées comp., invariantes par
\struck{\ill{}} $H$, sont en corr.\ biunivoque avec les orbites de $T$ dans
\struck{$X \smallsetminus X^{H}$} $X \smallsetminus X^{H}$ [dont les orbites
de \struck{\ill{}} $H^{+}$ dans $X \smallsetminus X^{H}$ définissent une
telle métrique]
\note{la première partie de la ligne, « $H^{+} = H \cap G^{+}$, $G^{+}$.
$X^{H} = X^{H^{+}}$, et $H^{+}$ », est en interligne au-dessus d'un
$H^{+}$ biffé ; le « $G^{+}$. » reste douteux}

\quad \struck{\ill{}} 2) Si $H = \{1, \sigma\}$, $\sigma$ une inversion, les
métriques invariantes sous $H$ sont en corr.\ biunivoque avec les éléments de
\[
  (X \smallsetminus X^{H})/H = (X \smallsetminus X^{\sigma})/\sigma
\]
\marginal{en marge droite, en oblique, le long d'un trait vertical qui court
sur toute la proposition : à vérifier. (\emph{NB} Il faudrait arriver à
traduire en termes d'alg.\ linéaire, l'action d'un autom.\ \struck{\ill{}}
de $G$ sur une métrique sphérique\,…)}

\page{37}
\note{pagination de l'auteur : 12}

\emph{Récapitulation}

\struck{\ill{}} \add{$X$} surface conforme (sans bord),
\struck{on peut lui} \add{\uncertain{On} \uncertain{peut}} définir
\add{\uncertain{compatible} \uncertain{avec} \uncertain{sa} \uncertain{str.}
\uncertain{conforme}} sur $X$ une métrique riemannienne canonique $Q_{X}$,
sauf dans les deux cas suivants :

a) $X \simeq \mathbb{C}$ (i.e.\ $X$ droite affine sur un corps
\add{\uncertain{top.}}/$\mathbb{C}$ isom.\ à $\mathbb{C}$), alors on peut
définir seulement une classe de métriques riemanniennes compatibles, mod.\
multiplication par un scalaire. Ce sont les métriques stables par
translations, qui sont aussi quasi-inv.\ par tt autom.\ conforme de $X$.
Elles sont \emph{plates} (donc à courbure nulle).

b) $X \simeq \mathbb{P}^{1}_{\mathbb{C}}$ (i.e.\ $X$ droite projective sur
un corps top.\ \uncertain{isom.}\ à $\mathbb{C}$, i.e.\ $X$ sphère
conforme). Alors il n'existe pas de métriques riemanniennes sur $X$
quasi-invariantes sous $G = \mathrm{Aut}\, X$, ni même sous
$\mathrm{Aut}^{+}(X)$. \struck{Pour} Il y a partout sur $X$ un \struck{\ill{}}
\add{ens.}\ $R(X)$ de métriques \struck{riemanniennes} \add{\ill{}}
riemanniennes (dites “sphériques normalisées”) avec $\int_{X} dx = 1$, qui
forment une orbite sous $G$. Si $H \subset G$ est un sous-groupe, il existe
un $Q \in R(X)$ stable par $H$ ssi $H$ est rel.\ cpt, et $Q$ est unique ssi
$X^{H} = \emptyset$. [Pour le cas $X^{H} \neq \emptyset$, cf.\ propositions
précédentes qui déterminent $R(X)^{H}$ dans tous les cas.]
\marginal{en marge gauche, en oblique : \emph{NB} $R(X)$ est une
\uncertain{variété} de \uncertain{dim} 3}

\page{38}
\note{pagination de l'auteur : 13}

[\emph{NB} \uncertain{Pour} \uncertain{passer} \ill{} au cas où $X$ est
\uncertain{associé} à une \emph{conique}, et $H$ est le groupe de ses
\uncertain{automorphismes}. \ill{} \uncertain{cette} \uncertain{conique},
qui est fini — la condition $X^{H} = \emptyset$ signifie aussi qu'il n'y a
pas de points de la conique fixés par $H$\,…]

Dans tous les cas où $Q_{X}$ est défini, $Q_{X}$ est hyperbolique (i.e.\ à
courbure $< 0$) ssi $X$ est hyperbolique, et alors $Q_{X}$ est de courbure
\add{constante} égale à $-1$. D'autre part, $Q_{X}$ est \emph{plate}
\struck{dans le cas} (courbure nulle) ssi
$\widetilde{X} \simeq \mathbb{C}$ [y inclus le cas a) ci-dessus, où on a
\uncertain{une} \uncertain{classe} de $Q_{X}$ mod.\ cte]. Le cas
\struck{restant} est celui du plan projectif réel, conforme,
\struck{pour lequel} \struck{\ill{}} \add{auquel} $Q_{X}$ est une métrique
homogène, de courbure cte \struck{égale à} $> 0$.

Donc chaque fois que $Q_{X}$ est défini, \uncertain{ou} \uncertain{défini}
mod.\ cte, il est à courbure \uncertain{constante}, \struck{\ill{}} Il en est
encore de $=$ de $Q_{X,H}$\note{le signe « $=$ » tient lieu de « même »} quand $X \simeq \mathbb{P}^{1}_{\mathbb{C}}$ et
$H \subset \mathrm{Aut}_{\mathrm{conf}}\, X$ est compact et tel que
$X^{H} = \emptyset$, on trouve \struck{\ill{}} des métriques
\uncertain{sph.}\ \uncertain{normalisées} \uncertain{dans} \ill{}, de
courbure constante $> 0$.
\marginal{en marge droite, en oblique, de haut en bas : \emph{NB} Pour
\uncertain{tt} $X$ 0-connexe \uncertain{conforme}, on a dans
l'\uncertain{ens.}\ \struck{\ill{}} $\mathrm{Riem}(X)$ des métriques
riemanniennes comp.\ avec la str.\ conforme une orbite canonique sous
$\mathrm{Aut}_{\mathrm{conf}}(X)$, qui se réduit à 1 pt dans tous les cas,
sauf les cas a), b) de la p.\ 12), où dans le cas a)), \uncertain{on}
\uncertain{trouve} une demi-droite dans $\mathrm{Riem}(X)$, et dans le cas
b), une variété de dim 3 (\struck{quelle} telle que $Q, Q' \in R(X)$,
$Q \neq Q'$ implique $Q'$ et $Q$ \emph{non} prop.). Notion des métriques
riemanniennes distinguées (a un sens aussi dans le cas des
\uncertain{variétés}\,?)}
\note{« la p.\ 12) » renvoie à sa pagination : la page 37, où sont les cas
a) et b)}

\page{39}
\note{pagination de l'auteur : 14}

\emph{Cas des surfaces conformes à bord} (pas néc.\ compactes)

On utilise l'équivalence de catégories
\[
  \text{surfaces conformes à bord} \;\approx\;
  \begin{array}{l}
    \text{surfaces holomorphes sans bord,} \\
    \text{munies d'une anti-involution}
  \end{array}
\]
\[
  X \longmapsto (X' = X^{\mathrm{or}}, \sigma)
\]
\[
  X = X'/\sigma \longleftarrow (X', \sigma)
\]
On dit que $X$ est hyperbolique ssi $X^{\mathrm{or}}$ est hyperbolique
(i.e.\ ses comp.\ connexes le sont).
Quand $X$ est sans bord, on a
$\widetilde{X} \simeq \widetilde{X^{\mathrm{or}}}$ (car $X^{\mathrm{or}}$ est
un rev.\ étale de $X$) donc $X$ est hyperbolique au sens des réflexions
précédentes, ssi il l'est en ce nouveau sens.
\note{l'exposant, écrit « or » dont le $r$ ressemble à un $2$, est rendu
$X^{\mathrm{or}}$ pour tout le lot : le revêtement double, surface
holomorphe sans bord munie de l'anti-involution $\sigma$}

Si $X$ est compacte, comme on a \add{dès $X^{\mathrm{or}}$ compacte}
\[
  \chi(X^{\mathrm{or}}) = 2\chi(X)
\]
et que $\chi(X^{\mathrm{or}}) < 0 \iff X^{\mathrm{or}}$ hyperbolique, on
trouve
\[
  X \text{ hyperbolique} \iff \chi(X) < 0
\]
P.\ ex.\ si $X$ est holomorphe 0-connexe, de type $(g, \nu)$ cela signifie
que $2g - 2 + \nu \geq 1$ i.e.\ $2g + \nu \geq 3$

Mais se garder de la confusion entre cette condition et celle que
$\operatorname{Int} X$ soit hyperbolique, ce qui est toujours vrai sauf si
$\nu = 0$, $g \in \{0, 1\}$\,…

\page{40}
\note{pagination de l'auteur : 15}

Si $X$ est une surface conforme hyperbolique, alors la métrique canonique
hyperbolique sur $X^{\mathrm{or}}$ se descend en une métrique hyperbolique
sur $X$, pour laquelle le bord est géodésique.

Plus généralement, chaque fois que $X^{\mathrm{or}}$ est muni d'une métrique
canonique, invariante par automorphismes conformes [ou seulement par
antiinvolution, ce qui est le cas aussi si
$X^{\mathrm{or}} \simeq \mathbb{C}$] alors celle-ci se descend à $X$, à bord
géodésique.

Quand $X$ est sans bord, \struck{\ill{}} alors \struck{la} les métriques
obtenues sur $X$ est celle déjà envisagée — \add{il y a}
\uncertain{unicité} dans tous les cas, sauf
$X \simeq \mathbb{P}^{1}_{\mathbb{C}}$. On fera attention cependant que
lorsque $X$ est le plan projectif réel, la deuxième prescription ne
\uncertain{fournit} plus de métrique — elle le fait si on exige de la
métrique sur $X^{\mathrm{or}}$ soit dans $R(X^{\mathrm{or}})$, que la
métrique sur $X^{\mathrm{or}}$ soit \emph{stable} par $\sigma$ (ce qui la
détermine \struck{\ill{}} sans ambiguïté). \struck{Dans}, dans le cas
$X \simeq \mathbb{C}$, $X^{\mathrm{or}} \simeq \mathbb{C} \amalg \mathbb{C}$,
la prescription donne un choix de métriques (mod.\ cte), la deuxième
\note{la page s'arrête au milieu de la phrase}

\end{document}
